\documentclass[12pt]{article}

% Idioma e codificação
\usepackage[brazil]{babel}
\usepackage[utf8]{inputenc}
\usepackage[T1]{fontenc}

% Matemática e listas
\usepackage{amsmath, amssymb}
\usepackage{enumitem}

% Layout e cores
\usepackage{geometry}
\usepackage{graphicx}
\usepackage{xcolor}

% Diagramas e Portas Lógicas
\usepackage{tikz}
\usetikzlibrary{shapes.gates.logic.US, positioning, arrows.meta, calc, fit, shadows.blur}

% Cores institucionais
\definecolor{maincolor}{RGB}{0, 79, 151}
\definecolor{meuazul}{RGB}{0, 79, 151}

% Cabeçalho e rodapé
\usepackage{fancyhdr}

% Blocos
\usepackage[most]{tcolorbox}
\newtcolorbox{block}[1]{%
  colback=blue!5,
  colframe=maincolor,
  coltitle=white,
  title=#1,
  fonttitle=\bfseries,
  boxrule=1.5pt,
  arc=3mm,
  left=6pt,
  right=6pt,
  top=6pt,
  bottom=6pt
}

\newtcolorbox{resposta}[1]{%
  colback=white,
  colframe=meuazul!30,
  arc=1mm,
  height=#1,
  width=\textwidth,
  boxrule=0.8pt
}

% Margens
\geometry{margin=2.5cm}

% Fancyhdr
\pagestyle{fancy}
\fancyhf{}

\renewcommand{\headrulewidth}{2pt}
\renewcommand{\footrulewidth}{2pt}

\renewcommand{\headrule}{\hbox to\headwidth{%
  \color{meuazul}\leaders\hrule height \headrulewidth\hfill}}

\renewcommand{\footrule}{\hbox to\headwidth{%
  \color{meuazul}\leaders\hrule height \footrulewidth\hfill}}

\setlength{\headheight}{52pt}
\addtolength{\topmargin}{-20pt}

\fancyhead[L]{%
  \includegraphics[height=1.2cm]{logo.png}
}

\fancyhead[R]{%
  \raggedleft
  \textcolor{meuazul}{%
    \small
    \textbf{Lista de Exercícios 4}\\
    Me.~Sergio Souza Novak\\
    \textit{Desenvolvimento de Software para Dispositivos Móveis}
  }
}

\fancyfoot[C]{\textcolor{meuazul}{\thepage}}

% =========================
\begin{document}
% =========================

\vspace*{0.2cm}

\begin{center}
  {\Large \textbf{Exercício de Implementação --- Jetpack Compose}}\\[0.2cm]
  \textit{Instruções: Implemente o exercício a seguir utilizando o Android Studio com Jetpack Compose.}\\[0.15cm]
  \textbf{Entrega:} O aluno deverá entregar \textbf{apenas um único arquivo Kotlin} (\texttt{MainActivity.kt})
  contendo todo o código-fonte, acompanhado de uma \textbf{pasta com as imagens de recurso} utilizadas no projeto (ex: fotos dos avatares).
\end{center}

\vspace{0.4cm}

% ── Exercício ───────────────────────────────────────────────────────────────
\textbf{Contexto:} Você foi contratado para desenvolver o protótipo de um aplicativo de videoconferências.
Sua tarefa é implementar a \textbf{tela de reunião em andamento} usando Jetpack Compose no Android Studio.

\vspace{0.3cm}
\textbf{Requisitos obrigatórios:}

\begin{enumerate}[label=\textbf{\arabic*.}]
  \item Criar um novo projeto no Android Studio com o template \textbf{Empty Activity} (Jetpack Compose).
  \item \textbf{Package name:} \texttt{com.seunome.meetingapp}
  \item A tela deve conter os seguintes elementos, conforme o wireframe abaixo:
  \begin{itemize}
    \item \textbf{Barra superior} com os nomes dos participantes da reunião (ex: ``Ethan, Alison, Daniela'') e a quantidade de pessoas (ex: ``8 people'').
    \item \textbf{Área central} com a foto (ou avatar circular) do participante em destaque e seu nome (ex: ``Ethan Martinez'').
    \item \textbf{Barra de controles inferior} com 3 botões de ação dispostos em \texttt{Row}:
      \begin{itemize}
        \item \textbf{Mute} (microfone) --- ícone com rótulo.
        \item \textbf{Raise} (levantar a mão) --- ícone com rótulo.
        \item \textbf{Audio} (alto-falante) --- ícone com rótulo.
      \end{itemize}
    \item \textbf{Botão de encerrar chamada} (vermelho, centralizado) abaixo dos controles.
  \end{itemize}
  \item Usar \texttt{data class} para representar um \textbf{Participante} (nome, foto).
  \item Criar componentes \texttt{@Composable} separados para: avatar do participante, botão de controle e barra de controles.
  \item Implementar \textbf{Preview} com \texttt{@Preview} para visualizar a tela sem executar o app.
\end{enumerate}

\vspace{0.4cm}
\textbf{Wireframe de referência:}

\vspace{0.3cm}
\begin{center}
\begin{tikzpicture}[
  scale=0.85, transform shape,
  % Estilos
  tela/.style={draw, rounded corners=12pt, minimum width=7cm, minimum height=12cm, fill=black!85, thick, draw=gray!50},
  barra/.style={draw=none, fill=black!70, minimum height=1.2cm, minimum width=6cm, rounded corners=4pt, align=center},
  avatar/.style={draw=gray!40, fill=gray!30, circle, minimum size=2.5cm, thick},
  btnctrl/.style={draw=gray!40, fill=gray!20, circle, minimum size=1.2cm, align=center},
  btnend/.style={draw=red!70!black, fill=red!65, rounded corners=14pt, minimum width=3.5cm, minimum height=0.85cm, align=center},
  lbl/.style={font=\tiny, text=white},
  lbldark/.style={font=\tiny, text=gray!60},
]

  % Corpo do telefone
  \node[tela] (phone) {};

  % --- Barra superior ---
  \node[barra, anchor=north, yshift=-0.3cm, font=\scriptsize\bfseries, text=white] (header) at (phone.north) {Ethan, Alison, Daniela};
  \node[lbldark, below=0.05cm of header] (pcount) {8 people};

  % --- Mini avatares dos outros participantes ---
  \node[draw=gray!40, fill=gray!25, circle, minimum size=0.9cm, thick] (p1) at ([yshift=2.8cm, xshift=-1.5cm]phone.center) {};
  \node[font=\tiny, text=gray!60] at (p1) {A};
  \node[lbl, below=0.08cm of p1, font=\tiny] {Alison};

  \node[draw=gray!40, fill=gray!25, circle, minimum size=0.9cm, thick] (p2) at ([yshift=2.8cm]phone.center) {};
  \node[font=\tiny, text=gray!60] at (p2) {D};
  \node[lbl, below=0.08cm of p2, font=\tiny] {Daniela};

  \node[draw=gray!40, fill=gray!25, circle, minimum size=0.9cm, thick] (p3) at ([yshift=2.8cm, xshift=1.5cm]phone.center) {};
  \node[font=\tiny, text=gray!60] at (p3) {L};
  \node[lbl, below=0.08cm of p3, font=\tiny] {Lucas};

  % --- Avatar central (participante em destaque) ---
  \node[avatar] (av) at ([yshift=0.5cm]phone.center) {};
  \node[font=\tiny, text=gray!50] at (av) {Foto};
  \node[lbl, below=0.3cm of av, font=\small\bfseries] (nome) {Ethan Martinez};

  % --- Botões de controle ---
  \node[btnctrl] (mute) at ([yshift=-3cm, xshift=-1.8cm]phone.center) {};
  \node[font=\tiny, text=gray!70] at (mute) {$\diagup$};
  \node[lbl, below=0.15cm of mute] {Mute};

  \node[btnctrl] (raise) at ([yshift=-3cm]phone.center) {};
  \node[font=\tiny, text=gray!70] at (raise) {$\uparrow$};
  \node[lbl, below=0.15cm of raise] {Raise};

  \node[btnctrl] (audio) at ([yshift=-3cm, xshift=1.8cm]phone.center) {};
  \node[font=\tiny, text=gray!70] at (audio) {$))$};
  \node[lbl, below=0.15cm of audio] {Audio};

  % --- Botão encerrar ---
  \node[btnend, font=\small\bfseries, text=white] (hangup) at ([yshift=-4.8cm]phone.center) {Encerrar};

  % --- Anotações laterais ---
  \draw[dashed, gray!50, -{Stealth[length=3pt]}] ([xshift=3.8cm]header.east) -- (header.east)
    node[pos=0, right, font=\tiny, text=gray!70, align=left] {Row \textbraceleft\\ \quad Text(participantes)\\ \quad Text(contagem)\\ \textbraceright};

  \draw[dashed, gray!50, -{Stealth[length=3pt]}] ([xshift=3.8cm]av.east) -- (av.east)
    node[pos=0, right, font=\tiny, text=gray!70, align=left] {Column \textbraceleft\\ \quad Image(avatar)\\ \quad Text(nome)\\ \textbraceright};

  \draw[dashed, gray!50, -{Stealth[length=3pt]}] ([xshift=3.8cm]raise.east) -- (audio.east)
    node[pos=0, right, font=\tiny, text=gray!70, align=left] {Row \textbraceleft\\ \quad BotaoControle(Mute)\\ \quad BotaoControle(Raise)\\ \quad BotaoControle(Audio)\\ \textbraceright};

  \draw[dashed, gray!50, -{Stealth[length=3pt]}] ([xshift=3.8cm]hangup.east) -- (hangup.east)
    node[pos=0, right, font=\tiny, text=gray!70, align=left] {Button(\\ \quad color = Vermelho,\\ \quad onClick = \textbraceleft...\textbraceright\\ )};

\end{tikzpicture}
\end{center}

\vspace{0.3cm}
\textbf{Dicas de componentes Compose:}
\begin{itemize}
  \item Use \texttt{Column}, \texttt{Row} e \texttt{Spacer} para organizar o layout.
  \item Use \texttt{Image} com \texttt{Modifier.clip(CircleShape)} para o avatar circular.
  \item Use \texttt{IconButton} ou \texttt{Button} para os controles.
  \item Use \texttt{MaterialTheme.colorScheme} para cores do tema.
  \item Fundo escuro: \texttt{Modifier.background(Color(0xFF1C1C1E))}.
\end{itemize}


\vspace{0.5cm}



\vspace{0.6cm}

\begin{center}
  \textit{``Boa atividade! Att. Me. Prof. Sergio Souza Novak.''}
\end{center}

\end{document}
% =========================
